\kbvcontinued ಪದಗಳು ಸೇರಿ ಸಮಸ್ತಪದವಾಗುವುದು. ಈ ಉದಾಹರಣೆಯಲ್ಲಿ “ತಾವರೆಕೊಳ” ಎಂಬುದು ಸಮಸ್ತಪದ. ಈ ಸಮಸ್ತಪದವೇ ಸಮಾಸ ಪದ.\par

\indent ಪೂರ್ವ ಪದ + ಉತ್ತರ ಪದ = ಸಮಸ್ತ ಪದ (ಸಮಾಸ)\par

\indent ಮೇಲಿನ ಉದಾಹರಣೆಯಲ್ಲಿ ಪೂರ್ವಪದವು ಷಷ್ಠೀ ವಿಭಕ್ತಿ ರೂಪದಲ್ಲಿದೆ; ಆದರೆ ಸಮಸ್ತಪದದಲ್ಲಿ ಈ ರೂಪ ಉಳಿದಿಲ್ಲ. ಆದುದರಿಂದ ಸಮಾಸವಾದಾಗ ವಿಭಕ್ತಿ ಪ್ರತ್ಯಯ ಲೋಪವಾದಂತಾಯಿತು.\par

\begin{kbvrule}
\noindent\kbvbullet\hspace{0.5em}ಸಮಸ್ತಪದವನ್ನು ಪೂರ್ವ ಮತ್ತು ಉತ್ತರಪದಗಳಾಗಿ ಬಿಡಿಸಿ ಬರೆಯುವ ಕ್ರಮಕ್ಕೆ ವಿಗ್ರಹ ವಾಕ್ಯ ಎಂದು ಹೆಸರು.
\end{kbvrule}

\vspace{3pt plus 2pt}\noindent\hspace*{1.2em}\kbvlabel{ಉದಾ}\par\vspace{1pt}

\begin{kbvderiv}
\vspace{4pt plus 3pt minus 1pt}
\begin{center}\begin{adjustbox}{max width=\textwidth}
{\renewcommand{\arraystretch}{1.22}\setlength{\tabcolsep}{4pt}
\begin{tabular}{l c l}
ಸಮಸ್ತಪದ & {\color{kbvgrey}$=$} & ಪೂರ್ವಪದ + ಉತ್ತರಪದ \\
ಮುಳ್ಳುಹಾಸಿಗೆ & {\color{kbvgrey}$=$} & ಮುಳ್ಳಿನ + ಹಾಸಿಗೆ \\
ದೊಡ್ಡತಲೆ & {\color{kbvgrey}$=$} & ದೊಡ್ಡದಾದ + ತಲೆ \\
ಹೊಸಗನ್ನಡ & {\color{kbvgrey}$=$} & ಹೊಸದು + ಕನ್ನಡ \\
\end{tabular}}
\end{adjustbox}\end{center}
\vspace{4pt plus 3pt minus 1pt}
\end{kbvderiv}

\begin{kbvrule}
\noindent\kbvbullet\hspace{0.5em}ಸಮಾಸಗಳಲ್ಲಿ ಹಲವು ಪ್ರಕಾರಗಳುಂಟು. ಅವುಗಳು ಯಾವುವೆಂಬುದನ್ನು ತಿಳಿಯೋಣ.
\end{kbvrule}

\vspace{4pt plus 4pt minus 1pt}\noindent{\kbvsecfont\fontsize{11pt}{14.5pt}\selectfont\color{kbvaccentdark} (೧)\hspace{0.45em}ತತ್ಪುರುಷ ಸಮಾಸ :}\hspace{0.5em}ಎರಡು ನಾಮಪದಗಳು ಸೇರಿ ಸಮಾಸವಾದಾಗ ಉತ್ತರ ಪದದ ಅರ್ಥವೇ ಪ್ರಧಾನವಾಗಿ ಉಳ್ಳ ಸಮಾಸವು ತತ್ಪುರುಷ ಸಮಾಸವೆನಿಸುವುದು. ಉದಾಹರಣೆಗೆ, “ತಲೆನೋವು” ಎಂಬ ಸಮಸ್ತಪದವನ್ನು ವಿಗ್ರಹವಾಕ್ಯ ಮಾಡಿನೋಡೋಣ.\par

\indent ತಲೆಯಲ್ಲಿ + ನೋವು = ತಲೆನೋವು ಇಲ್ಲಿ ಪೂರ್ವಪದ ಮತ್ತು ಉತ್ತರಪದಗಳಲ್ಲಿ ಉತ್ತರಪದದ ಅರ್ಥವೇ ಪ್ರಧಾನವಾಗಿದೆ; ಅಂದರೆ, ದೇಹದಲ್ಲಾಗುವ ನೋವು ಪ್ರಧಾನವೇ ಹೊರತು. ದೇಹದ ಯಾವ ಭಾಗದಲ್ಲಿ ನೋವಾಗುತ್ತಿರುವುದೆಂಬುದು ಅಷ್ಟು ಪ್ರಧಾನವಲ್ಲ, ಆದುದರಿಂದ, ಇಲ್ಲಿನ ಸಮಾಸವು ತತ್ಪುರುಷ ಸಮಾಸವಾಗಿರುವುದು.\par

\begin{kbvrule}
\noindent\kbvbullet\hspace{0.5em}ತತ್ಪುರುಷ ಸಮಾಸದಲ್ಲಿ ಪೂರ್ವಪದವು ತೃತೀಯಾ ವಿಭಕ್ತಿಯಿಂದ ಮೊದಲ್ಗೊಂಡು ಸಪ್ತಮೀ ವಿಭಕ್ತಿಯವರೆಗೆ ಯಾವುದಾದರೊಂದು
\end{kbvrule}
